% GENERATED FILE. Edit canonical lessons, then run npm run generate:books.
% canonical-source-hash: fnv1a64:128be4c136db8049
% canonical-lessons: GU-C391-gadbad, GU-C391-avyavastha, GU-C391-bhangvu, GU-C391-veravu, GU-C391-bhusavu

\chapter{When Things Go Wrong}
\label{ch:gu-when-things-go-wrong}
\hlchaptermodality{391}

\begin{chapteropening}
\textbf{By the end of this chapter:} \emph{I can warn what will go wrong if plans fall apart.}

Complete the last lesson of chapter 391: I can warn what will go wrong if plans fall apart.
\end{chapteropening}

\section[\texorpdfstring{\gu{ગડબડ} (gaḍbaḍ)}{ગડબડ (gaḍbaḍ)}]{\gu{ગડબડ} (gaḍbaḍ) — a mix-up}

\label{lesson:GU-C391-gadbad}



\begin{quote}
Before the new one: say the Gujarati for to grow plants, then the Gujarati for to be produced.
\end{quote}

\subsection*{You'll want to know: \gu{ગડબડ}}
\textbf{\gu{ગડબડ}} — \emph{gaḍbaḍ} — \textquotedblleft{}a mix-up\textquotedblright{}.

\begin{quote}
\gu{જો} \gu{ટિકિટની} \gu{તારીખમાં} \gu{ગડબડ} \gu{હશે} \gu{તો} \gu{મુસાફરી} \gu{રદ} \gu{થશે}.
\end{quote}

\begin{quote}
\emph{jo ṭikiṭnī tārīkhmā̃ gaḍbaḍ haśe to musāpharī rad thaśe.}
\end{quote}

\begin{quote}
If there is a mix-up over the date on the ticket, the trip will be cancelled.
\end{quote}

A feminine noun.

\begin{cousinweb}
\textbf{\gu{ગડબડ}} (\emph{gaḍbaḍ}) CDIAL lists the Gujarati word, \textquotedblleft{}confused; confusion, disorder\textquotedblright{}, under a reconstructed \emph{gaḍa-baḍa} (CDIAL 3974), a jingle of sounds for a muddle.
\end{cousinweb}

\begin{grammarlens}[title={What you've built}]
One more word for saying what happens if.
\end{grammarlens}

\subsection*{Your turn}
\begin{itemize}
\raggedright
  \item \emph{Say it:} \emph{gaḍbaḍ}
  \item \emph{Say it:} \emph{gaḍbaḍ}, once more
  \item \emph{Say it:} \emph{jo ṭikiṭnī tārīkhmā̃ gaḍbaḍ haśe to musāpharī rad thaśe.} — If there is a mix-up over the date on the ticket, the trip will be cancelled.
  \item \emph{Recall:} say the Gujarati for to grow plants, then the Gujarati for to be produced, then say \emph{gaḍbaḍ} again
\end{itemize}

\begin{tcolorbox}[breakable,colback=teal!4,colframe=teal!35!black,title={Before you move on}]
What does \gu{ગડબડ} mean? (\textquotedblleft{}A mix-up\textquotedblright{}.) Say it once more.
\end{tcolorbox}
\section[\texorpdfstring{\gu{અવ્યવસ્થા} (avyavasthā)}{અવ્યવસ્થા (avyavasthā)}]{\gu{અવ્યવસ્થા} (avyavasthā) — disorder}

\label{lesson:GU-C391-avyavastha}



\begin{quote}
Before the new one: say the Gujarati for to be produced, then the Gujarati for a mix-up.
\end{quote}

\subsection*{You'll want to know: \gu{અવ્યવસ્થા}}
\textbf{\gu{અવ્યવસ્થા}} — \emph{avyavasthā} — \textquotedblleft{}disorder\textquotedblright{}.

\begin{quote}
\gu{જો} \gu{સ્ટેશન} \gu{પર} \gu{ભીડ} \gu{વધશે} \gu{તો} \gu{અવ્યવસ્થા} \gu{ફેલાશે}.
\end{quote}

\begin{quote}
\emph{jo sṭeśan par bhīḍ vadhśe to avyavasthā phelāśe.}
\end{quote}

\begin{quote}
If the crowd at the station grows, disorder will spread.
\end{quote}

A feminine noun.

\begin{cousinweb}
\textbf{\gu{અવ્યવસ્થા}} (\emph{avyavasthā}) is a tatsama, borrowed from Sanskrit \emph{avyavasthā}, \textquotedblleft{}irregularity, want of order\textquotedblright{} (Monier-Williams).
\end{cousinweb}

\begin{grammarlens}[title={What you've built}]
One more word for saying what happens if.
\end{grammarlens}

\subsection*{Your turn}
\begin{itemize}
\raggedright
  \item \emph{Say it:} \emph{avyavasthā}
  \item \emph{Say it:} \emph{avyavasthā}, once more
  \item \emph{Say it:} \emph{jo sṭeśan par bhīḍ vadhśe to avyavasthā phelāśe.} — If the crowd at the station grows, disorder will spread.
  \item \emph{Recall:} say the Gujarati for to be produced, then the Gujarati for a mix-up, then say \emph{avyavasthā} again
\end{itemize}

\begin{tcolorbox}[breakable,colback=teal!4,colframe=teal!35!black,title={Before you move on}]
What does \gu{અવ્યવસ્થા} mean? (\textquotedblleft{}Disorder\textquotedblright{}.) Say it once more.
\end{tcolorbox}
\section[\texorpdfstring{\gu{ભાંગવું} (bhā̃gvũ)}{ભાંગવું (bhā̃gvũ)}]{\gu{ભાંગવું} (bhā̃gvũ) — to shatter}

\label{lesson:GU-C391-bhangvu}



\begin{quote}
Before the new one: say the Gujarati for a mix-up, then the Gujarati for disorder.
\end{quote}

\subsection*{You'll want to know: \gu{ભાંગવું}}
\textbf{\gu{ભાંગવું}} — \emph{bhā̃gvũ} — \textquotedblleft{}to shatter\textquotedblright{}.

\begin{quote}
\gu{જો} \gu{તમે} \gu{ભારે} \gu{થેલી} \gu{ખુરશી} \gu{પર} \gu{મૂકશો} \gu{તો} \gu{ખુરશી} \gu{ભાંગશે}.
\end{quote}

\begin{quote}
\emph{jo tame bhāre thelī khurśī par mūkśo to khurśī bhā̃gśe.}
\end{quote}

\begin{quote}
If you put the heavy bag on the chair, the chair will break.
\end{quote}

Past (it …): \gu{તે} \gu{જૂનો} \gu{પુલ} \gu{ભાંગ્યો}, \emph{te jūno pul bhā̃gyo}.

\begin{cousinweb}
\textbf{\gu{ભાંગવું}} (\emph{bhā̃gvũ}) CDIAL lists the Gujarati verb, \textquotedblleft{}to break\textquotedblright{}, under Sanskrit \emph{bhaṅgá}, \textquotedblleft{}breaking\textquotedblright{} (CDIAL 9353).
\end{cousinweb}

\begin{grammarlens}[title={What you've built}]
One more word for saying what happens if.
\end{grammarlens}

\subsection*{Your turn}
\begin{itemize}
\raggedright
  \item \emph{Say it:} \emph{bhā̃gvũ}
  \item \emph{Say it:} \emph{bhā̃gvũ}, once more
  \item \emph{Say it:} \emph{jo tame bhāre thelī khurśī par mūkśo to khurśī bhā̃gśe.} — If you put the heavy bag on the chair, the chair will break.
  \item \emph{Recall:} say the Gujarati for a mix-up, then the Gujarati for disorder, then say \emph{bhā̃gvũ} again
\end{itemize}

\begin{tcolorbox}[breakable,colback=teal!4,colframe=teal!35!black,title={Before you move on}]
What does \gu{ભાંગવું} mean? (\textquotedblleft{}To shatter\textquotedblright{}.) Say it once more.
\end{tcolorbox}
\section[\texorpdfstring{\gu{વેરાવું} (verāvũ)}{વેરાવું (verāvũ)}]{\gu{વેરાવું} (verāvũ) — to get scattered}

\label{lesson:GU-C391-veravu}



\begin{quote}
Before the new one: say the Gujarati for disorder, then the Gujarati for to shatter.
\end{quote}

\subsection*{You'll want to know: \gu{વેરાવું}}
\textbf{\gu{વેરાવું}} — \emph{verāvũ} — \textquotedblleft{}to get scattered\textquotedblright{}.

\begin{quote}
\gu{જો} \gu{તમે} \gu{ડબ્બો} \gu{ઉતાવળમાં} \gu{ખોલશો} \gu{તો} \gu{બધી} \gu{ખાંડ} \gu{વેરાઈ} \gu{જશે}.
\end{quote}

\begin{quote}
\emph{jo tame ḍabbo utāvaḷmā̃ kholśo to badhī khā̃ḍ verāī jaśe.}
\end{quote}

\begin{quote}
If you open the tin in a hurry, the sugar will spill everywhere.
\end{quote}

Past (it …): \gu{તે} \gu{ચોખા} \gu{વેરાઈ} \gu{ગયા}, \emph{te cokhā verāī gayā}.

\begin{cousinweb}
\textbf{\gu{વેરાવું}} (\emph{verāvũ}) is the \emph{-āvũ} partner of \emph{vervũ}, \textquotedblleft{}to scatter\textquotedblright{}; CDIAL lists that verb under a reconstructed \emph{vikērayati}, \textquotedblleft{}scatters\textquotedblright{} (CDIAL 11636).
\end{cousinweb}

\begin{grammarlens}[title={What you've built}]
One more word for saying what happens if.
\end{grammarlens}

\subsection*{Your turn}
\begin{itemize}
\raggedright
  \item \emph{Say it:} \emph{verāvũ}
  \item \emph{Say it:} \emph{verāvũ}, once more
  \item \emph{Say it:} \emph{jo tame ḍabbo utāvaḷmā̃ kholśo to badhī khā̃ḍ verāī jaśe.} — If you open the tin in a hurry, the sugar will spill everywhere.
  \item \emph{Recall:} say the Gujarati for disorder, then the Gujarati for to shatter, then say \emph{verāvũ} again
\end{itemize}

\begin{tcolorbox}[breakable,colback=teal!4,colframe=teal!35!black,title={Before you move on}]
What does \gu{વેરાવું} mean? (\textquotedblleft{}To get scattered\textquotedblright{}.) Say it once more.
\end{tcolorbox}
\section[\texorpdfstring{\gu{ભૂંસાવું} (bhū̃sāvũ)}{ભૂંસાવું (bhū̃sāvũ)}]{\gu{ભૂંસાવું} (bhū̃sāvũ) — to get wiped out}

\label{lesson:GU-C391-bhusavu}



\begin{quote}
Before the new one: say the Gujarati for to shatter, then the Gujarati for to get scattered.
\end{quote}

\subsection*{You'll want to know: \gu{ભૂંસાવું}}
\textbf{\gu{ભૂંસાવું}} — \emph{bhū̃sāvũ} — \textquotedblleft{}to get wiped out\textquotedblright{}.

\begin{quote}
\gu{જો} \gu{વીજળી} \gu{જશે} \gu{તો} \gu{સ્ક્રીન} \gu{પરનું} \gu{બધું} \gu{કામ} \gu{ભૂંસાઈ} \gu{જશે}.
\end{quote}

\begin{quote}
\emph{jo vījḷī jaśe to skrīn parnũ badhũ kām bhū̃sāī jaśe.}
\end{quote}

\begin{quote}
If the power goes, all the work on the screen will be wiped out.
\end{quote}

Past (it …): \gu{તે} \gu{પાટિયા} \gu{પરનું} \gu{નામ} \gu{ભૂંસાયું}, \emph{te pāṭiyā parnũ nām bhū̃sāyũ}.

\begin{cousinweb}
\textbf{\gu{ભૂંસાવું}} (\emph{bhū̃sāvũ}) is the \emph{-āvũ} partner of \textbf{\gu{ભૂંસવું}} (\emph{bhū̃svũ}), \textquotedblleft{}to rub out\textquotedblright{}, whose own history is not settled.
\end{cousinweb}

\begin{grammarlens}[title={What you've built}]
One more word for saying what happens if.
\end{grammarlens}

\subsection*{Your turn}
\begin{itemize}
\raggedright
  \item \emph{Say it:} \emph{bhū̃sāvũ}
  \item \emph{Say it:} \emph{bhū̃sāvũ}, once more
  \item \emph{Say it:} \emph{jo vījḷī jaśe to skrīn parnũ badhũ kām bhū̃sāī jaśe.} — If the power goes, all the work on the screen will be wiped out.
  \item \emph{Recall:} say the Gujarati for to shatter, then the Gujarati for to get scattered, then say \emph{bhū̃sāvũ} again
\end{itemize}

\begin{tcolorbox}[breakable,colback=teal!4,colframe=teal!35!black,title={Before you move on}]
What does \gu{ભૂંસાવું} mean? (\textquotedblleft{}To get wiped out\textquotedblright{}.) Say it once more.
\end{tcolorbox}
